% Section 6.2: 研究创新点总结

在 6.1 节研究工作总结基础上，本节抽取本文区别于现有工作的核心创新点，并从方法设计、实验协议与工程实现三个角度加以归纳。

\subsection{训练侧与推理侧解耦的二因素析因设计}
以往工作常把「是否微调」与「是否 CoT」耦合在单一对照中，无法分离训练与推理各自的贡献。本文将两者严格解耦为独立因子，并引入「训练/推理 $k$ 匹配」合法性约束，把 $4 \times 3 = 12$ 个理论组合收敛为 10 组有效实验 H1--H10。该设计首次在同一口径下同时估计了：
\begin{itemize}
    \item 训练侧主效应（$H4/H7/H9 - H1$）；
    \item 推理侧主效应（$H2/H3 - H1$）；
    \item 四条训练 $\times$ 推理交互项；
    \item few-shot 边际在不同训练侧上的一致性。
\end{itemize}
这一析因结构是本文方法侧的主要增量。

\subsection{Few-shot CoT 在训练侧与推理侧的对称化}
本文把 few-shot 示例同时注入训练 prompt 与推理 prompt，共享同一评分函数与同一候选池规则，并通过 \texttt{seed = 2026 + problem\_idx * 1000 + sample\_idx} 使评测侧抽样完全可复现。该对称化直接支撑「训练与推理 $k$ 同步放大是否带来净增益」「$\mathrm{H9} \gtrsim \mathrm{H10}$ 是否成立」等 5.3 节设问，突破了以往 few-shot 仅作用于推理侧的常见做法。

\subsection{两阶段 CoT 流水线的工程约束}
为让 CoT 在 10 组评测中稳定产出可比较的 pass@$k$，本文在 4.4 节提出三项工程约束：
\begin{itemize}
    \item \textbf{推理段二级截断}：字符级 + token 级联合裁剪，防止 Stage-1 压缩 Stage-2 预算；
    \item \textbf{Stage-1 提前代码块裁断}：避免推理偷跑代码污染 Stage-2 上下文；
    \item \textbf{非法组合脚本级拒绝}：在评测脚本内硬编码「训练/推理 $k$ 匹配」，防止误操作回流到矩阵外。
\end{itemize}
这些约束共同构成本文结论可复现的工程保障。

\subsection{端到端可复现流水线}
从 Parquet 原始数据到 H1--H10 十份 JSON 结果，全流程由 \texttt{scripts/train\_eval.sh} 一键驱动，训练与评测脚本顶部的超参数、指针文件机制、结果命名规则均在仓库内显式固定。任何第三方只需下载基座权重并打开对应开关，即可完整重跑本文结论。